\subsection{Annäherung von links und von rechts}
\label{subsec:real_functions_one_sided_limits}
Das Verhalten kann davon abhängen, von welcher Seite wir uns nähern.
Deshalb unterscheiden wir \emph{einseitige Grenzwerte}
\EnglishTerm{one-sided limits}:
\[
  \lim_{x\to a^-}f(x)
  \quad\text{mit }x<a,
  \qquad
  \lim_{x\to a^+}f(x)
  \quad\text{mit }x>a.
\]
\begin{ExampleBox}{Ein Sprung hat zwei verschiedene Annäherungswerte}{real_functions_jump_limit}
  Eine idealisierte Schaltfunktion sei durch
  $s(x)=0$ für $x<0$ und $s(x)=1$ für $x\ge0$ gegeben.
  \par\smallskip
  \begin{minipage}[c]{0.47\linewidth}
    Beide Seiten liefern ein eindeutiges, aber unterschiedliches Ergebnis:
    \[
      \lim_{x\to0^-}s(x)=\Cloze{1cm}{0},
      \qquad
      \lim_{x\to0^+}s(x)=\Cloze{1cm}{1}.
    \]
    Der zweiseitige Grenzwert bei null existiert nicht.
    Auch eine Änderung von $s(0)$ würde daran nichts ändern.
  \end{minipage}\hfill
  \begin{minipage}[c]{0.49\linewidth}
    \centering
    \CachedTikzFigure{content/resources/chapters/real_functions/tikz/functions_jump.tex}
  \end{minipage}
\end{ExampleBox}
\begin{KeyPointBox}{Beide Seiten müssen zusammenpassen}{real_functions_two_sided_criterion}
  Wenn eine Annäherung von beiden Seiten möglich ist, gilt für $L\in\R$:
  \[
    \lim_{x\to a}f(x)=L
    \quad\Longleftrightarrow\quad
    \lim_{x\to a^-}f(x)=\lim_{x\to a^+}f(x)=L.
  \]
  Am Rand der Definitionsmenge betrachten wir nur die zulässige Seite,
  etwa $x\to0^+$ bei $\sqrt{x}$ auf $[0,\infty)$.
\end{KeyPointBox}

\subsection{Unbegrenzte Eingaben und unbegrenzte Funktionswerte}
\label{subsec:real_functions_infinite_limits}
Die Schreibweise mit $\infty$ kann zwei verschiedene Fragen ausdrücken.
Man muss unterscheiden, \emph{was} unbegrenzt wird.
\par\noindent\begin{minipage}{\linewidth}
\begin{DefinitionBox}{Grenzverhalten im Unendlichen}{real_functions_limit_at_infinity}
  Bei $x\to+\infty$ wächst die Eingabe über jede feste Schranke,
  bei $x\to-\infty$ wird sie beliebig stark negativ.

  Die Aussage $\lim_{x\to+\infty}f(x)=L$ bedeutet:
  Die Funktionswerte liegen und bleiben beliebig nahe bei $L$,
  sobald $x$ genügend groß ist. Für $x\to-\infty$ gilt dies entsprechend.
\end{DefinitionBox}
\begin{ExampleBox}[before skip=0.3\baselineskip]{Die Eingabe wächst -- die Ausgabe nähert sich null}{real_functions_reciprocal_infinity}
  Für große positive Eingaben vergleichen wir die Werte von $f(x)=1/x$:
  \[
    \begin{array}{c|rrrr}
      x&10&100&1000&10000\\ \hline
      1/x&0{,}1&0{,}01&0{,}001&0{,}0001
    \end{array}
    \qquad
    \lim_{x\to+\infty}\frac1x=\Cloze{1cm}{0}.
  \]
  Der Wert null wird nie angenommen. Ein Grenzwert muss also kein
  Funktionswert sein. Auch für $x\to-\infty$ gilt $1/x\to0$.
\end{ExampleBox}
\end{minipage}


\par\noindent\begin{minipage}{\linewidth}
\begin{DefinitionBox}{Uneigentliches Grenzverhalten}{real_functions_divergence}
  Die Schreibweise $f(x)\to+\infty$ bedeutet, dass $f(x)$ schließlich
  jede feste obere Schranke überschreitet und darüber bleibt.
  Entsprechend bedeutet $f(x)\to-\infty$ ein dauerhaftes Unterschreiten
  jeder festen unteren Schranke. Man spricht von einem
  \emph{uneigentlichen Grenzwert} \EnglishTerm{infinite limit}.
\end{DefinitionBox}
\begin{ExampleBox}[before skip=0.3\baselineskip]{Eine endliche Stelle -- unbegrenzte Werte}{real_functions_reciprocal_zero}
  Bei $f(x)=1/x$ nähern wir uns nun nicht dem Unendlichen, sondern der
  ausgeschlossenen Stelle null:
  \[
    \lim_{x\to0^-}\frac1x=\Cloze{2cm}{-\infty},
    \qquad
    \lim_{x\to0^+}\frac1x=\Cloze{2cm}{+\infty}.
  \]
  Es gibt bei null weder einen gemeinsamen endlichen noch einen
  gemeinsamen uneigentlichen Grenzwert.
  \begin{center}
    \CachedTikzFigure{content/resources/chapters/real_functions/tikz/functions_reciprocal.tex}
  \end{center}
\end{ExampleBox}
\end{minipage}

\begin{RemarkBox}
  $\infty$ ist keine reelle Zahl. Insbesondere sind
  $1/0$, $\infty-\infty$ und $\infty/\infty$ keine zulässigen
  gewöhnlichen Rechnungen.
\end{RemarkBox}
